\documentclass{article}
\usepackage{lmodern}
\usepackage{amsmath}
\usepackage{amssymb}
\usepackage[margin=1in]{geometry}

\title{The TeXbook Mathematical Typesetting Reference Suite}
\author{Donald E. Knuth (Specimens from Chapters 16--19)}
\date{}

\begin{document}
\maketitle

\section*{Chapter 16: Typing Math Formulas}

\subsection*{16.1 Basic Variables, Digits, and Signs}
The simplest formula is a single letter or number:
$x$, $2$, $x + y - z$, $(x + y)/(x - y) = 1$.

\subsection*{16.2 Greek Letters}
Lowercase Greek:
$\alpha, \beta, \gamma, \delta, \epsilon, \varepsilon, \zeta, \eta, \theta, \vartheta, \iota, \kappa, \lambda, \mu, \nu, \xi, \pi, \varpi, \rho, \varrho, \sigma, \varsigma, \tau, \upsilon, \phi, \varphi, \chi, \psi, \omega$.

Uppercase Greek:
$\Gamma, \Delta, \Theta, \Lambda, \Xi, \Pi, \Sigma, \Upsilon, \Phi, \Psi, \Omega$.

\subsection*{16.3 Subscripts and Superscripts}
Basic scripts:
$x^2$, $x_i$, $x_i^2$, $x_{i,j}^2$, $2^{2^{2^x}}$, $A_{i,j,k}^{2\prime}$.

Nested script styles:
$y_{x_2}$, $a_{b_c}$, $e^{-x^2}$, $2^{n+1}$, $x_1^{\prime\prime\prime}$.

\subsection*{16.4 Square Roots and Radicals}
Radicals of varying heights:
$\sqrt{x}$, $\sqrt{x+y}$, $\sqrt{x^2 + y^2}$, $\sqrt{n+1}$, $\sqrt{\frac{x}{y}}$, $\sqrt{x_1 + \sqrt{x_2 + \sqrt{x_3}}}$.

\subsection*{16.5 Fractions}
Inline and displayed fractions:
$\frac{1}{2}$, $\frac{x+y}{x-y}$, $\frac{a}{b+c}$.

Nested fractions:
\begin{equation}
a_0 + \frac{1}{a_1 + \frac{1}{a_2 + \frac{1}{a_3}}}
\end{equation}

\subsection*{16.6 Math Accents}
Standard mathematical accents:
$\hat{a}$, $\check{a}$, $\tilde{a}$, $\acute{a}$, $\grave{a}$, $\dot{a}$, $\ddot{a}$, $\breve{a}$, $\bar{a}$, $\vec{a}$.

Accents over uppercase letters:
$\hat{A}$, $\hat{L}$, $\hat{H}$, $\vec{B}$, $\bar{X}$.

Wide accents:
$\widehat{x+y}$, $\widetilde{A+B}$.

\section*{Chapter 17: More about Math}

\subsection*{17.1 Delimiters and Fences}
Various delimiter shapes:
$(x)$, $[x]$, $\{x\}$, $|x|$, $\|x\|$, $\langle x \rangle$, $\lfloor x \rfloor$, $\lceil x \rceil$.

Automatically scaled delimiters:
\begin{equation}
\left( \frac{x^2 + 1}{x^2 - 1} \right), \quad
\left[ 1 + \left( \frac{a}{b} \right)^2 \right], \quad
\left| \frac{\sqrt{x} - 1}{\sqrt{x} + 1} \right|
\end{equation}

Manual delimiter sizing:
$\bigl( x \bigr) \quad \Bigl( x \Bigr) \quad \biggl( x \biggr) \quad \Biggl( x \Biggr)$.

\subsection*{17.2 Large Operators}
Summations and products:
\begin{equation}
\sum_{i=1}^n i = \frac{n(n+1)}{2}, \quad \prod_{j=1}^m \left( 1 - \frac{1}{p_j^s} \right)
\end{equation}

Integrals and contour integrals:
\begin{equation}
\int_0^\infty e^{-x^2} \, dx = \frac{\sqrt{\pi}}{2}, \quad \oint_C \mathbf{F} \cdot d\mathbf{r} = 0
\end{equation}

Unions and intersections:
\begin{equation}
\bigcup_{k=1}^\infty A_k, \quad \bigcap_{k=1}^\infty B_k
\end{equation}

\section*{Chapter 18: Fine Points of Mathematics Typing}

\subsection*{18.1 Spacing Between Atom Classes}
Binary operator vs. unary sign:
$x + y$ vs. $+y$, and $a - b$ vs. $-b$.

Differential spacing with thin space:
\begin{equation}
\iint_D f(x, y) \, dx \, dy
\end{equation}

Factorials and punctuation:
$n! / (k! \, (n-k)!)$, and $f(x) : X \to Y$, and $(x_1, x_2, \dots, x_n)$.

\subsection*{18.2 Ellipses}
Horizontal, centered, and vertical dots:
$x_1 + x_2 + \dots + x_n$, and $x_1 x_2 \cdots x_n$.

\section*{Chapter 19: Displayed Equations}

\subsection*{19.1 Numbered and Centered Displays}
\begin{equation}
e^{i\pi} + 1 = 0
\end{equation}

\begin{equation}
[\hat{x}_j, \hat{p}_k] = i\hbar\delta_{jk}, \quad [\hat{L}_x, \hat{L}_y] = i\hbar\hat{L}_z
\end{equation}

\begin{equation}
a^\dagger |n\rangle = \sqrt{n+1} |n+1\rangle, \quad a |n\rangle = \sqrt{n} |n-1\rangle
\end{equation}

\subsection*{19.2 Custom Equation Tags and Unnumbered Displays}
Custom named tag via \verb|\tag{...}|:
\begin{equation}
\nabla \times \mathbf{E} = -\frac{\partial \mathbf{B}}{\partial t} \tag{Faraday}
\end{equation}

Unparenthesized tag via \verb|\tag*{...}|:
\begin{equation}
\nabla \cdot \mathbf{B} = 0 \tag*{[Gauss]}
\end{equation}

Unnumbered display equation via \verb|\begin{equation*}|:
\begin{equation*}
\lim_{n\to\infty} \left( 1 + \frac{1}{n} \right)^n = e
\end{equation*}

\subsection*{19.3 Multiline Alignments}
Aligned multiline equation system via \verb|\begin{align*}|:
\begin{align*}
(x + y)^2 &= (x + y)(x + y) \\
&= x^2 + 2xy + y^2
\end{align*}

\section*{Chapter 26: Summary of Math Mode and Spacing Primitives}

\subsection*{26.1 Phantoms and Metric Reservations}
Tensor indices with phantom vertical and horizontal spacing:
\begin{equation}
R^{\mu\phantom{\nu}\rho}_{\phantom{\mu}\nu\phantom{\rho}\sigma}
\end{equation}

Horizontal and vertical phantoms:
\begin{equation}
\left( \frac{1}{\vphantom{y^2}x} \right) \quad \text{vs.} \quad \left( \frac{1}{y^2} \right)
\end{equation}

\subsection*{26.2 Smashes and Overlaps}
Smashed expression without vertical bounding box expansion:
\begin{equation}
\sqrt{\smash{x^2 + y^2}} + \sqrt{a + b}
\end{equation}

\section*{Appendix G: Generating Boxes from Formulas (The 22 Rules)}

\subsection*{G.1 Axis Centering and Vcenter (Rules 8 \& 12)}
\begin{equation}
A + \vcenter{\hbox{\large$\frac{x}{y}$}} = B
\end{equation}

\subsection*{G.2 Overlines and Underlines (Rules 9 \& 10)}
\begin{equation}
\overline{\overline{x + y}} \quad \text{and} \quad \underline{\underline{a - b}}
\end{equation}

\subsection*{G.3 Radical Clearances and Root Indices (Rule 11)}
\begin{equation}
\sqrt{\sqrt{x^4 + 1}} + \sqrt[\beta]{\alpha^2 + \gamma^2} = \sqrt{\frac{a^2 + b^2}{c^2 + d^2}}
\end{equation}

\subsection*{G.4 Big Operators and Limits (Rules 13 \& 13a)}
\begin{equation}
\sum_{k=1}^\infty \frac{1}{k^2} = \frac{\pi^2}{6}, \quad \prod_{p \text{ prime}} \left(1 - \frac{1}{p^s}\right)^{-1} = \zeta(s)
\end{equation}

\subsection*{G.5 Generalized Fractions and Binomials (Rules 15a--15e)}
\begin{equation}
\binom{2n}{n} = \frac{(2n)!}{(n!)^2}, \quad \frac{1 + \frac{x}{y}}{1 - \frac{x}{y}} = \frac{y + x}{y - x}
\end{equation}

\subsection*{G.6 Subscript and Superscript Combinations (Rules 18a--18f)}
\begin{equation}
R^{\mu\nu}_{\rho\sigma} + g^{\mu\alpha} \Gamma_{\alpha\beta}^\nu x^\beta = 0, \quad x_{i_1 i_2}^{j_1 j_2}
\end{equation}

\subsection*{G.7 Symmetric Delimiter Sizing (Rule 19)}
\begin{equation}
\left( \sum_{k=1}^n \frac{1}{k} \right) \left[ \prod_{j=1}^m \left( 1 - \frac{1}{p_j} \right) \right] \left\{ \int_0^1 \frac{dx}{\sqrt{1 - x^2}} \right\}
\end{equation}

\end{document}


